\documentclass[11pt]{article}

\usepackage{amsmath, amssymb}
\usepackage{geometry}
\usepackage{graphicx}
\usepackage{booktabs}

\geometry{margin=1in}

\title{Portfolio Risk and Scenario Analysis Engine}
\author{Robin Basso}
\date{\today}

\begin{document}

\maketitle

\section{Introduction}

This project implements a portfolio risk engine designed to evaluate the distribution of outcomes and downside risk under both normal market conditions and stress scenarios.

The primary objective is to provide decision-relevant metrics for portfolio managers and family offices, focusing on capital preservation and downside risk.

\section{Portfolio Model}

We consider a portfolio composed of $N$ assets with weights:

\[
\mathbf{w} = (w_1, w_2, ..., w_N), \quad \sum_{i=1}^N w_i = 1
\]

Let $\mathbf{r}_t$ denote the vector of asset returns at time $t$:

\[
\mathbf{r}_t = (r_{t,1}, r_{t,2}, ..., r_{t,N})
\]

The portfolio return at time $t$ is:

\[
R_t = \mathbf{w}^\top \mathbf{r}_t
\]

The portfolio value evolves as:

\[
V_t = V_0 \prod_{s=1}^{t} (1 + R_s)
\]

\section{Stochastic Returns Model}

Asset returns are modeled as a multivariate Gaussian process:

\[
\mathbf{r}_t \sim \mathcal{N}(\boldsymbol{\mu}, \boldsymbol{\Sigma})
\]

where:
\begin{itemize}
\item $\boldsymbol{\mu}$ is the vector of expected returns
\item $\boldsymbol{\Sigma}$ is the covariance matrix
\end{itemize}

Correlation between assets is captured through $\boldsymbol{\Sigma}$, ensuring realistic joint behavior.

\section{Monte Carlo Simulation}

We simulate $M$ independent scenarios of length $T$:

\[
\{\mathbf{r}_t^{(m)}\}_{t=1}^T, \quad m = 1, ..., M
\]

For each simulation:
\begin{enumerate}
\item Generate asset returns from the multivariate distribution
\item Compute portfolio returns
\item Construct the portfolio value path
\end{enumerate}

This yields a distribution of terminal values:

\[
\{V_T^{(m)}\}_{m=1}^M
\]

\section{Risk Metrics}

\subsection{Value at Risk (VaR)}

Value at Risk at confidence level $\alpha$ is defined as:

\[
\text{VaR}_\alpha = V_0 - Q_{1-\alpha}(V_T)
\]

where $Q_{1-\alpha}$ is the $(1-\alpha)$ quantile of the terminal value distribution.

\subsection{Conditional Value at Risk (CVaR)}

Conditional VaR (Expected Shortfall) measures the expected loss beyond VaR:

\[
\text{CVaR}_\alpha = \mathbb{E}[V_0 - V_T \mid V_T \leq Q_{1-\alpha}(V_T)]
\]

\subsection{Maximum Drawdown}

For a given path $\{V_t\}$, define the running maximum:

\[
M_t = \max_{s \leq t} V_s
\]

The drawdown is:

\[
D_t = \frac{M_t - V_t}{M_t}
\]

Maximum drawdown is:

\[
\text{MDD} = \max_t D_t
\]

\section{Scenario Analysis}

Stress scenarios modify the return process to reflect adverse market conditions.

\subsection{Equity Crash}

An instantaneous shock applied at $t=0$:

\[
r_{0,\text{equity}} \rightarrow r_{0,\text{equity}} - \delta
\]

with $\delta \approx 20\%$.

\subsection{Volatility Spike}

Increased uncertainty is modeled by scaling returns:

\[
\mathbf{r}_t \rightarrow \lambda \mathbf{r}_t, \quad \lambda > 1
\]

\subsection{Combined Stress}

A combination of shocks:

\begin{itemize}
\item Equity crash
\item Increased volatility
\end{itemize}

\section{Interpretation}

The model allows answering key portfolio questions:

\begin{itemize}
\item What is the worst expected loss at a given confidence level?
\item How severe are losses in extreme scenarios?
\item How does the portfolio behave under market stress?
\end{itemize}

\section{Implementation Notes}

The implementation follows a modular architecture:

\begin{itemize}
\item \texttt{Portfolio}: handles portfolio aggregation
\item \texttt{ReturnsModel}: generates stochastic returns
\item \texttt{MonteCarloEngine}: simulates paths
\item \texttt{RiskMetrics}: computes risk measures
\item \texttt{Scenario}: applies stress transformations
\end{itemize}

The design emphasizes:
\begin{itemize}
\item clarity
\item extensibility
\item separation of concerns
\end{itemize}

\section{Conclusion}

This project provides a practical framework for portfolio risk analysis, bridging quantitative modeling and decision-making tools for investment professionals.

\end{document}